%==============================================================
%  ito_foundations.tex
%  hasklib — mathematical companion / derivations
%
%  Foundations: probability space, Brownian motion, the SDE
%  framework, Ito's lemma, quadratic variation.
%
%  This file is the base every other derivation file references
%  (gbm.tex, ou.tex, abm.tex, cir.tex, schemes.tex all cite results here),
%  exactly as every process class inherits from Process.hpp.
%
%  Scope discipline: include a result ONLY if a piece of hasklib
%  depends on it. No Levy processes, no jumps, nothing the library
%  does not touch.
%==============================================================
\documentclass[11pt]{article}

%---------------- packages ----------------
\usepackage[utf8]{inputenc}
\usepackage[T1]{fontenc}
\usepackage{amsmath, amssymb, amsthm}
\usepackage{mathtools}
\usepackage[a4paper, margin=1in]{geometry}
\usepackage{xcolor}
\usepackage{hyperref}
\usepackage{cleveref}
\usepackage{bm}
\usepackage{harvard}   % Harvard author-date referencing (with agsm.bst)
\usepackage{zref-clever}   % lets gbm.tex/ou.tex/etc. \zcref{} into this file's theorems

%---------------- notation macros ----------------
% Probability / expectation
\newcommand{\Prob}{\mathbb{P}}          % real-world measure
\newcommand{\Qmeas}{\mathbb{Q}}         % risk-neutral measure
\newcommand{\E}{\mathbb{E}}             % expectation
\newcommand{\Var}{\operatorname{Var}}  % variance
\newcommand{\Cov}{\operatorname{Cov}}  % covariance
\newcommand{\filt}{\mathcal{F}}        % sigma-algebra / filtration
\newcommand{\Om}{\Omega}               % sample space

% Calculus
\newcommand{\dd}{\mathrm{d}}           % upright d for differentials
\newcommand{\dW}{\dd W_t}              % Brownian increment
\newcommand{\dt}{\dd t}                % time increment
\newcommand{\R}{\mathbb{R}}

% Common process objects
\newcommand{\drift}{a}                 % drift coefficient a(t,x)
\newcommand{\diff}{b}                  % diffusion coefficient b(t,x)

%---------------- theorem environments ----------------
\theoremstyle{plain}
\newtheorem{theorem}{Theorem}
\newtheorem{lemma}{Lemma}
\newtheorem{proposition}{Proposition}

\theoremstyle{definition}
\newtheorem{definition}{Definition}
\newtheorem{assumption}{Assumption}

\theoremstyle{remark}
\newtheorem{remark}{Remark}

% A visible marker for gaps you still need to fill in.
\newcommand{\TODO}[1]{\medskip\noindent\textbf{[\textcolor{red}{TO DERIVE}: #1]}\medskip}

%==============================================================
\title{Foundations: It\^o Calculus and the SDE Framework\\
       \large \texttt{hasklib} documentation}
\author{Hubert Kasinski}
\date{\today}
%==============================================================

\begin{document}
\maketitle

\begin{abstract}
This note collects the stochastic-calculus results that the
\texttt{stochastic/} layer of \texttt{hasklib} relies on: the
probability-space setup, Brownian motion, the general It\^o SDE
$\dd X_t = \drift(t,X_t)\dt + \diff(t,X_t)\dW$ (encoded by the
\texttt{StochasticProcess} interface), It\^o's lemma, and quadratic
variation. Downstream files (\texttt{gbm.tex}, \texttt{ou.tex},
\texttt{abm.tex}, \texttt{cir.tex}, \texttt{schemes.tex}) cite the results established
here. Standard measure-theoretic probability is assumed; see
\citeasnoun{shreve2004} or \citeasnoun{williams1991} for construction.
\end{abstract}

%--------------------------------------------------------------
\section{Probability space and Brownian motion}
%--------------------------------------------------------------

We work on a filtered probability space
$(\Om, \filt, \{\filt_t\}_{t\ge 0}, \Prob)$ satisfying the usual
conditions. The filtration $\{\filt_t\}$ represents the information
available up to time $t$.

\begin{definition}[Brownian motion]
A standard Brownian motion $\{W_t\}_{t\ge 0}$ is a process with
\begin{enumerate}
  \item $W_0 = 0$ almost surely;
  \item Independent increments: for $s < t < u $, $W_u - W_t$ is
        independent of $W_t - W_s$;
  \item Gaussian increments: $W_t - W_s \sim \mathcal{N}(0,\, t-s)$;
  \item Continuous sample paths.
\end{enumerate}
\end{definition}

\begin{remark}
In code, property (3) is what \texttt{NormalRng::draw()} supplies:
a step of length $\Delta t$ uses $W_{t+\Delta t}-W_t = \sqrt{\Delta t}\,Z$
with $Z\sim\mathcal N(0,1)$. This substitution is the single bridge
between the continuous theory here and the discrete simulation code.
\\
\end{remark}


Although the paths are continuous, they are extremely rough. Fix $t$ and
consider the difference quotient over a step $h>0$. By the Gaussian-increment
property, $W_{t+h}-W_t \sim \mathcal{N}(0,h)$, equivalently $\sqrt{h}\,Z$ with
$Z\sim\mathcal{N}(0,1)$, so
\[
   \frac{W_{t+h}-W_t}{h} \;\overset{d}{=}\; \frac{1}{\sqrt{h}}\,Z,
\]
whose standard deviation $1/\sqrt{h}$ diverges as $h\downarrow 0$. Informally,
the path moves by order $\sqrt{h}$ over a time step of order $h$, so its
instantaneous slope is of order $1/\sqrt{h}\to\infty$ and the quotient
converges to no finite limit. This heuristic is made rigorous by the
Paley--Wiener--Zygmund theorem \cite{pwz1933}: with probability one,
$t\mapsto W_t$ is nowhere differentiable.

Equivalently, the path has infinite first-order (total) variation on every
interval, so integrals $\int f\,\dd W_s$ cannot be defined pathwise as
Riemann--Stieltjes integrals. Its quadratic variation, however, is
finite, and it is precisely this finite second-order variation that the It\^o integral and It\^o's lemma are built on.

%--------------------------------------------------------------
\section{Quadratic variation}\label{sec:qv}
%--------------------------------------------------------------

The property that distinguishes stochastic from ordinary calculus is
that Brownian motion accumulates quadratic variation at a
nonzero rate.

\begin{proposition}[Quadratic variation of Brownian motion]\label{prop:qv}
For a standard Brownian motion, the quadratic variation over $[0,T]$
satisfies
\[
   \langle W\rangle_T \;=\; \lim_{\|\Pi\|\to 0}\sum_{k} \big(W_{t_{k+1}}-W_{t_k}\big)^2
        \;=\; T
   \qquad \text{in } L^2,
\]
where $\Pi$ is a partition of $[0,T]$. Informally this is written
\[
   (\dW)^2 = \dt.
\]
\end{proposition}

\begin{proof}
Write the partition as $\Pi=\{0=t_0<t_1<\dots<t_n=T\}$ and abbreviate
\[
  \Delta t_k = t_{k+1}-t_k, \qquad
  \Delta W_k = W_{t_{k+1}}-W_{t_k}, \qquad
  S_\Pi = \sum_{k=0}^{n-1}\big(\Delta W_k\big)^2 .
\]
``Convergence in $L^2$'' means the mean-square error vanishes:
$\E\big[(S_\Pi-T)^2\big]\to0$ as the mesh $\|\Pi\|:=\max_k\Delta t_k\to0$.
We use only two facts about the increments, both directly from the definition
of Brownian motion --- each $\Delta W_k\sim\mathcal N(0,\Delta t_k)$, and
increments over disjoint intervals are independent.

\medskip
\noindent\emph{Step 1:}\ The mean is exactly $T$.
\\Since $\E[\Delta W_k]=0$, the second moment is the variance:
\[
  \E\big[(\Delta W_k)^2\big]=\Var(\Delta W_k)=\Delta t_k .
\]
Summing over $k$,
\[
  \E[S_\Pi]=\sum_{k}\E\big[(\Delta W_k)^2\big]
           =\sum_k \Delta t_k = T ,
\]
which holds for every partition, not merely in the limit. Because the
mean is exactly $T$, the mean-square error is the variance:
\[
  \E\big[(S_\Pi-T)^2\big]=\Var(S_\Pi) .
\]
So the whole statement reduces to showing $\Var(S_\Pi)\to0$.

\medskip
\noindent\emph{Step 2}\; Variance of a single term.
\\
Normalise the increment as $\Delta W_k=\sqrt{\Delta t_k}\,Z$ with
$Z\sim\mathcal N(0,1)$, so $(\Delta W_k)^2=\Delta t_k\,Z^2$ and
\[
  \Var\big((\Delta W_k)^2\big)=(\Delta t_k)^2\,\Var(Z^2).
\]
For a standard normal, $\E[Z^2]=1$ and $\E[Z^4]=3$, hence
$\Var(Z^2)=\E[Z^4]-\big(\E[Z^2]\big)^2=3-1=2$. (The fourth moment is one
Gaussian integration by parts: with the density $\phi$ satisfying
$\phi'(z)=-z\,\phi(z)$, integrating $\int z^4\phi\,\dd z$ by parts gives
$\E[Z^4]=3\,\E[Z^2]=3$.) Therefore
\[
  \Var\big((\Delta W_k)^2\big)=2\,(\Delta t_k)^2 .
\]

\medskip
\noindent\emph{Step 3}\; Variance of the sum, then let the mesh shrink.
\\
The increments $\Delta W_k$ are independent, so the squared increments
$(\Delta W_k)^2$ are independent too, and the variance of a sum of
independent variables is the sum of the variances:
\[
  \Var(S_\Pi)=\sum_k \Var\big((\Delta W_k)^2\big)=2\sum_k (\Delta t_k)^2 .
\]
Bound one factor of $\Delta t_k$ in each term by the mesh:
\[
  \sum_k (\Delta t_k)^2 \le \Big(\max_k \Delta t_k\Big)\sum_k \Delta t_k
                        = \|\Pi\|\,T .
\]
Combining,
\[
  \E\big[(S_\Pi-T)^2\big]=\Var(S_\Pi)\le 2\,\|\Pi\|\,T
  \;\xrightarrow[\;\|\Pi\|\to0\;]{}\;0 ,
\]
which is exactly the claimed $L^2$ convergence.
\end{proof}

\begin{remark}
The mean computation ($\E[S_\Pi]=T$) is what makes $(\dW)^2=\dt$
plausible; the variance computation is what makes it usable.
Per step the squared increment $(\Delta W_k)^2$ still fluctuates --- its
standard deviation $\sqrt2\,\Delta t_k$ is the same order as its mean --- but
those fluctuations are independent across steps, so summing averages them out
at rate $\|\Pi\|$. Inside a sum (equivalently, an integral) the random
$(\Delta W_k)^2$ may therefore be replaced by the deterministic $\Delta t_k$.
That replacement is the entire content of the shorthand $(\dW)^2=\dt$ used in
the next section.
\end{remark}

%--------------------------------------------------------------
\section{The It\^o stochastic differential equation}
%--------------------------------------------------------------

\begin{definition}[It\^o SDE]
An It\^o process $X_t$ solves the stochastic differential equation
\begin{equation}\label{eq:sde}
   \dd X_t \;=\; \drift(t,X_t)\,\dt \;+\; \diff(t,X_t)\,\dW,
\end{equation}
interpreted as the integral equation
\[
   X_t = X_0 + \int_0^t \drift(s,X_s)\,\dd s
              + \int_0^t \diff(s,X_s)\,\dd W_s,
\]
where the second integral is an It\^o integral. We call $\drift$ the
drift and $\diff$ the diffusion coefficient.
\end{definition}

\begin{remark}[Correspondence with the code]
Equation~\eqref{eq:sde} is exactly what the \texttt{StochasticProcess}
base class encodes: \texttt{drift(t,x)} returns $\drift(t,x)$ and
\texttt{diffusion(t,x)} returns $\diff(t,x)$. Every concrete process
(GBM, OU, CIR) is a choice of these two functions.
\end{remark}

\begin{assumption}[Existence and uniqueness]
We assume $\drift$ and $\diff$ satisfy the standard Lipschitz and
linear-growth conditions guaranteeing a unique strong solution to
\eqref{eq:sde}; see \citeasnoun{oksendal2003}, Thm.~5.2.1. \emph{(Cited, not proved ---
the library never depends on the proof, only on the conclusion.)}
\end{assumption}

%--------------------------------------------------------------
\section{It\^o's lemma}
%--------------------------------------------------------------

It\^o's lemma is the chain rule for functions of an It\^o process, and
it is the engine behind every closed-form solution used in
\texttt{hasklib} (the GBM exact sampler, the OU solution).

\begin{theorem}[It\^o's lemma, one dimension]\label{thm:ito}
Let $X_t$ solve \eqref{eq:sde} and let $\varphi(t,x)\in C^{1,2}$. Then
\begin{equation}\label{eq:ito}
   \dd \varphi(t,X_t) =
     \Big( \frac{\partial\varphi}{\partial t} + \drift\, \frac{\partial\varphi}{\partial x} + \tfrac12 \diff^2 \frac{\partial^2\varphi}{\partial x^2} \Big)\dt
     \;+\; \diff\, \frac{\partial\varphi}{\partial x} \,\dW,
\end{equation}
\end{theorem}

\begin{proof}[Heuristic derivation]
This is the standard second-order Taylor argument. It is not a proof in the
strict sense --- a rigorous version requires the construction of the It\^o
integral \cite[Ch.~4]{shreve2004}, which is outside the scope of this note ---
but it exposes exactly where the extra term comes from, and it is all the
library needs.

The one idea is this: because $\dd X$ carries a Brownian term of size
$\sqrt{\dt}$, a contribution that ordinary calculus discards as ``second
order'' is in fact first order and must be kept.

\medskip
\noindent\emph{Orders of magnitude.}\;
Over a step of length $\dt$ the increment $\dW$ has typical size $\sqrt{\dt}$
(it is $\mathcal N(0,\dt)$). Ranking the elementary products by their power of
$\dt$:
\[
  \dt = O(\dt), \qquad
  \dW = O(\dt^{1/2}), \qquad
  (\dW)^2 = O(\dt), \qquad
  \dt\,\dW = O(\dt^{3/2}), \qquad
  (\dt)^2 = O(\dt^2).
\]
Discard everything of order smaller than $\dt$, and use \cref{prop:qv} to
replace $(\dW)^2$ by its in-sum limit $\dt$. This gives the \emph{It\^o
multiplication table}
\begin{equation}\label{eq:itotable}
  (\dW)^2=\dt, \qquad \dt\,\dW=0, \qquad (\dt)^2=0 .
\end{equation}
The last two entries are not literally zero; they are negligible next to $\dt$
and vanish in the mesh limit for exactly the reason the fluctuations did in
\cref{sec:qv}.

\medskip
\noindent\emph{Taylor expansion.}\;
Expand $\varphi(t+\dt,\,X_t+\dd X)$ about $(t,X_t)$ to second order:
\[
  \dd \varphi = \frac{\partial\varphi}{\partial t}\,\dt + \frac{\partial\varphi}{\partial x}\,\dd X
        + \tfrac12 \frac{\partial^2\varphi}{\partial t^2}\,(\dt)^2
        + \frac{\partial^2\varphi}{\partial t\partial x}\,\dt\,\dd X
        + \tfrac12 \frac{\partial^2\varphi}{\partial x^2}\,(\dd X)^2
        + \cdots
\]
Substitute the SDE $\dd X = \drift\,\dt + \diff\,\dW$ and square:
\[
  (\dd X)^2 = \drift^2(\dt)^2 + 2\,\drift\diff\,\dt\,\dW + \diff^2(\dW)^2 .
\]
Apply \eqref{eq:itotable}. In $(\dd X)^2$ the first two terms die and only
$\diff^2(\dW)^2=\diff^2\,\dt$ survives:
\[
  (\dd X)^2 = \diff^2\,\dt .
\]
The two remaining second-order terms of the expansion vanish for the same
reason:
\[
  \tfrac12 \frac{\partial^2\varphi}{\partial t^2}\,(\dt)^2 = 0, \qquad
  \frac{\partial^2\varphi}{\partial t\partial x}\,\dt\,\dd X
     = \frac{\partial^2\varphi}{\partial t\partial x}\big(\drift\,(\dt)^2 + \diff\,\dt\,\dW\big) = 0 .
\]
What remains is
\[
  \dd \varphi = \frac{\partial\varphi}{\partial t}\,\dt + \frac{\partial\varphi}{\partial x}\,\dd X + \tfrac12 \frac{\partial^2\varphi}{\partial x^2}\,\diff^2\,\dt .
\]
Substitute $\dd X = \drift\,\dt + \diff\,\dW$ once more and collect the $\dt$
and $\dW$ parts:
\[
  \dd \varphi = \Big( \frac{\partial\varphi}{\partial t} + \drift\, \frac{\partial\varphi}{\partial x} + \tfrac12 \diff^2 \frac{\partial^2\varphi}{\partial x^2} \Big)\dt
        + \diff\, \frac{\partial\varphi}{\partial x} \,\dW ,
\]
which is \eqref{eq:ito}.
\end{proof}

\begin{remark}
Had $X$ been a smooth (finite-variation) path, $(\dd X)^2$ would have been
genuinely $O((\dt)^2)$ and dropped out, leaving the ordinary chain rule
$\dd \varphi = \frac{\partial\varphi}{\partial t}\,\dt + \frac{\partial\varphi}{\partial x}\,\dd X$. The whole gap between stochastic and ordinary
calculus is the single surviving term $\tfrac12\diff^2 \frac{\partial^2\varphi}{\partial x^2}\,\dt$, and it
survives only because $(\dW)^2=\dt$ rather than $0$ --- i.e.\ because Brownian
motion has nonzero quadratic variation. Every closed-form result in
\texttt{hasklib} that ``corrects a drift by half a variance'' (the
$-\tfrac12\sigma^2$ in the GBM log-solution, the It\^o term in the OU solution)
is exactly this one term.
\end{remark}

\begin{remark}
The single downstream consequence to keep in view: applying
\cref{thm:ito} with $f(t,x)=\ln x$ to geometric Brownian motion
linearises the SDE and yields the exact terminal law sampled by
\texttt{GBM::sample\_terminal}.
\end{remark}

%--------------------------------------------------------------
\section{What the following files derive}
%--------------------------------------------------------------

Each companion file starts from \eqref{eq:sde} with a specific
$(\drift,\diff)$ and uses \cref{thm:ito} to produce results the code
depends on:
\begin{itemize}
  \item \texttt{gbm.tex}: $\drift=\mu x,\ \diff=\sigma x$ $\Rightarrow$
        exact solution and moments (verified in
        \texttt{test\_stochastic.cpp}).
  \item \texttt{ou.tex}: $\drift=\kappa(\theta-x),\ \diff=\sigma$
        $\Rightarrow$ Gaussian solution, mean/variance, stationary law.
  \item \texttt{abm.tex}: $\drift=\mu,\ \diff=\sigma$ $\Rightarrow$
        exact solution by direct integration (no It\^o's lemma needed);
        the $\kappa=0$ degenerate case of \texttt{ou.tex}, plus a
        drift.
  \item \texttt{cir.tex}: $\drift=\kappa(\theta-x),\ \diff=\sigma\sqrt{x}$
        $\Rightarrow$ moments, Feller condition, motivation for full
        truncation.
  \item \texttt{schemes.tex}: discretisation of \eqref{eq:sde} ---
        Euler--Maruyama (strong order $0.5$) and Milstein
        (strong order $1.0$) via stochastic Taylor expansion
        \cite{kloeden1992}.
\end{itemize}

%--------------------------------------------------------------
\bibliographystyle{agsm}
\bibliography{refs}
%--------------------------------------------------------------

\end{document}